\chapter{Units, vectors and conservation}\label{ch:units}
Many apparent disagreements about EBU are disagreements about units. One
author writes $q$ for a rate, another for a finite quantity; both use the
same symbol in a correct-looking equation. The error appears only when the
duration differs from one. We will therefore identify quantity before
performing algebra.

\section{A dimensional dictionary}
Write $\X$ for the unit of physical stock and $T$ for time. In the active
Gaussian normalization, $V$ is dimensionless, as are its finite differences.
Calling that numerical accounting unit EBU does not make it a joule, a litre
or a currency unit. A later calibrated conversion must be declared rather
than inferred from the name.

\begin{center}
\begin{tabular}{P{.39\linewidth}P{.48\linewidth}}
\toprule Object & Unit and meaning\\\midrule
$x_i,x_i^*,\sigma_i,q$ & $\X$: stock, reference, scale, finite withdrawal\\
$j$ & $\X/T$: a rate\\
$\mu_i=\partial V/\partial x_i$ & $1/\X$: potential change per unit stock\\
$H_{ii}=1/\sigma_i^2$ & $1/\X^2$: curvature\\
$q\mu_i$, $q^2H_{ii}$ & dimensionless finite-account terms\\
\bottomrule
\end{tabular}
\end{center}

If a device operates at constant rate $j$ for duration $\Delta t$, its
quantity is $q=\Delta t\,j$. Thus $\delta=S q$ and
$\delta=\Delta t S j$ are equivalent descriptions, not two steps to apply
successively. The historical foundation uses rate notation
\cite{foundation}; this edition uses finite quantities in its active examples.

\section{Reading a matrix multiplication}
For a lossless chain $1\to2\to3$, let
\begin{equation}
 S=\begin{pmatrix}-1&0\\1&-1\\0&1\end{pmatrix},
 \qquad q=\begin{pmatrix}a\\b\end{pmatrix}.
\end{equation}
Then
\begin{equation}
 Sq=\begin{pmatrix}-a\\a-b\\b\end{pmatrix}.
\end{equation}
The first cell loses $a$, the middle receives $a$ and sends $b$, and the
last receives $b$. The middle coordinate is not a transport label that may
be omitted: it is a physical stock with its own predecessor and successor.

Every column of $S$ sums to zero. If $\mathbf1$ is the vector of ones,
$\mathbf1^T S=0$, so
\begin{equation}
 \mathbf1^T(x+Sq)=\mathbf1^Tx.
\end{equation}
This proves conservation for the declared lossless increments. It does not
prove nonnegativity: a request $a>x_1$ still conserves the algebraic total
while making the source negative. Conservation and admissibility must be
checked separately.

\section{Loss is a physical account, not a vanished number}
For a lossy transfer, a column may have $-1$ at the source and $\eta$ at the
destination. Its sum is $-(1-\eta)$. To close a broader physical account,
the missing quantity must enter a declared loss carrier or cross an explicit
boundary. For instance an added coordinate $\ell$ can receive
$(1-\eta)q$. That is a conservation declaration, not yet a valuation of
$\ell$ or permission to charge a second process cost.

The active first domain sets $\eta=1$. This does not dismiss loss. It avoids
claiming that an unresolved loss-aware runtime and valuation convention has
already been accepted. Part III returns to the old Definition 6.4 wording:
adding a loss coordinate changes what can honestly be said about a carrier
being absent from state.

\section{Changing the physical unit}
Suppose litres are replaced by millilitres. Multiply all stocks, references,
scales and finite quantities by $c=1000$. Each standardized deviation remains
unchanged:
\begin{equation}
 \frac{c x_i-c x_i^*}{c\sigma_i}
 =\frac{x_i-x_i^*}{\sigma_i}.
\end{equation}
Consequently $V$ and exact EBU are unchanged. The marginal becomes
$\mu_i/c$, while $q$ becomes $cq$; their product is unchanged. This is a
useful dimensional check on an implementation specification.

Changing $\sigma$ while keeping stocks fixed is not a unit conversion. It
changes the model's relative sensitivity. Treating those operations as the
same would hide a substantive calibration decision inside a formatting step.

\section*{Worked checks}
A rate of six units per hour over a quarter hour transfers $3/2$ units.
Over two hours it transfers twelve. A routine that uses six in both cases
has silently assumed a unit-duration interval.

In the chain matrix, take $a=3$, $b=2$ and predecessor $(5,1,0)$.
The declared endpoint is $(2,2,2)$, with total six before and after. Whether
the middle cell may send two while initially holding one depends on timing:
one may allow a declared inflow-funded simultaneous process, or require
outgoing commitments to be covered by frozen stock. The matrix alone cannot
settle that physical convention.

\section*{Chapter recap}
Vectors collect quantities; matrices express declared transformations. Units
and column sums expose important errors, but they do not certify the whole
physical process. Conservation, nonnegativity, timing and valuation each
need their own check.
